\documentclass[10pt,a4paper]{amsart}
\usepackage[margin=19mm]{geometry}
\usepackage{amsmath,amssymb,mathtools}
\usepackage[scr=rsfs,cal=euler]{mathalfa}
\usepackage{xfrac}
\usepackage[unicode,hidelinks,bookmarks=false]{hyperref}
\newcommand{\ie}{i.e.\ }
\newcommand{\cf}{cf.\ }
\newcommand{\resp}{respectively\ }
\newcommand{\Subsection}{\subsection}
\providecommand{\nobreakdash}{\nobreak\hbox{-}}
\setlength{\emergencystretch}{2em}
\multlinegap0pt
\usepackage{amssymb}
\usepackage{bm}
\usepackage{bbm}
\usepackage{enumerate}
\newtheorem{definition}{Definition}
\newtheorem{theorem}{Theorem}
\newcommand{\GL}{{\rm GL}}\newcommand{\SL}{{\rm SL}}
\newcommand{\bF}{\mathbbm{F}}\newcommand{\bP}{\mathbbm{P}}
\newcommand{\bZ}{\mathbbm{Z}}\newcommand{\bQ}{\mathbbm{Q}}
\title{Birational classification for algebraic tori}
\author{Akinari Hoshi, Aiichi Yamasaki}
\date{Source: arXiv:2112.02280v9 (CC BY 4.0)}
\begin{document}
\maketitle

\noindent\emph{Source and licence.} Birational classification for algebraic tori. Version arXiv:2112.02280v9 (CC BY 4.0), distributed by arXiv under the Creative Commons Attribution 4.0 International licence, http://creativecommons.org/licenses/by/4.0/, as recorded in the arXiv licence metadata for this exact version and shown on the abstract page https://arxiv.org/abs/2112.02280v9. Full licence notice: \texttt{licenses/arxiv-2112.02280-CC-BY-4.0.txt}.\par\smallskip

\noindent\textit{Editorial scope.} Counted unit: the article's theorem thmXYZ (source0.tex lines 1279--1325). The article's complete original proof of it is delivered with it (source0.tex lines 1327--1365, one \texttt{proof} environment of 39 lines). No mathematics is written, repaired or re-derived: the statement and the proof are the article's own lines, in the article's own notation, and no retained line is substituted or re-flowed.  Retained uncounted as the source-local context the proof needs: lines 678--680; lines 1112--1119; lines 1247--1274; lines 2022--2174.  Nothing was omitted from any retained span, so the package carries no omission record.  No retained line contains a citation, so the wrapper carries no bibliography and no bibliography entry.  No retained line contains a figure environment, an image-inclusion command or drawing code, so no image is delivered and the package declares no asset dependency.  Editorial formatting only, in addition: an AMS \texttt{amsart} preamble that restates the article's own declarations and macros used by the retained lines, namely the environments \texttt{definition}, \texttt{theorem} on separate counters, together with the standard packages those lines need.  The retained environments print in this document as Definition 1, Theorem 1 in order of appearance, whereas the article prints the same items with the numbering of its own full document.  The complete native source of the article as distributed in the arXiv e-print for this exact version is bundled unchanged as \texttt{sources/119254/source0.tex}.
\begin{align}
x_i^{\sigma}=\prod_{j=1}^n x_j^{a_{i,j}},\quad 1\leq i\leq n\label{acts}
\end{align}

\begin{definition}[The $G$-lattice $M_G$ of a finite subgroup $G$ of 
$\GL(n,\bZ)$]\label{d2.2} 
Let $G$ be a finite subgroup of $\GL(n,\bZ)$. 
{\it The $G$-lattice $M_G$ of rank $n$} 
is defined to be the $G$-lattice with a $\bZ$-basis $\{u_1,\ldots,u_n\}$ 
on which $G$ acts by $u_i^{\sigma}=\sum_{j=1}^n a_{i,j}u_j$ for any $
\sigma=[a_{i,j}]\in G$ (corresponding to the equation (\ref{acts})). 
\end{definition}

\begin{definition}\label{defXYZ}
We define the following subgroups of ${\rm Aut}(G)$ for 
$G\leq {\rm GL}(n\,\bZ)$: 
\begin{center}
${\rm Inn}(G)\leq X\leq Y\leq Z\leq {\rm Aut}(G)$,
\end{center}
\begin{align*}
X&={\rm Aut}_{\GL(n,\bZ)}(G)=
\{\sigma\in{\rm Aut}(G)\mid 
{\rm there}\ {\rm exists}\ u\in {\rm GL}(n,\bZ)\ {\rm such}\ {\rm that}\ 
u^{-1}gu=g^\sigma\ {\rm for}\ {\rm any}\ g\in G\}\\
&\simeq N_{\GL(n,\bZ)}(G)/Z_{\GL(n,\bZ)}(G),\\
Y&=\{\sigma\in{\rm Aut}(G)\mid [M_G]^{fl}=[M_{G^\sigma}]^{fl}\ {\rm as}\ 
\widetilde{H}\textrm{-{\rm lattices}}\ {\rm where}\ \widetilde{H}=\{(g,g^\sigma)\mid g\in G\}\simeq G\},\\
Z&=\{\sigma\in{\rm Aut}(G)\mid [M_H]^{fl}\sim [M_{H^\sigma}]^{fl}\ {\rm for}\ {\rm any}\ H\leq G\} 
\end{align*}
where 
${\rm Inn}(G)$ is the group of inner automorphisms on $G$, 
${\rm Aut}(G)$ is the group of automorphisms on $G$, 
$N_{\GL(n,\bZ)}(G)$ is the normalizer of $G$ in $\GL(n,\bZ)$ and 
$Z_{\GL(n,\bZ)}(G)$ is the centralizer of $G$ in $\GL(n,\bZ)$.
%Under the assumption $L=L^\prime$, i.e. 
%$\widetilde{H}\simeq G\simeq G^\sigma$, 
%
%In particular, if ${\rm Inn}(G)={\rm Aut}(G)$, then 
%${\rm (1)}\Leftrightarrow {\rm (2)}$ of Theorem \ref{thmain4} holds.
%
\end{definition}

\begin{theorem}[Main theorem 4]\label{thmain4}
Let $k$ be a field.  
Let $T_i$ and $T_j^\prime$ $(1\leq i,j\leq 152)$ 
be algebraic $k$-tori of dimension $4$ 
with the minimal splitting fields $L_i$ and $L^\prime_j$ 
and the character modules 
$\widehat{T}_i=M_G$ and $\widehat{T}^\prime_j=M_{G^\prime}$ 
which satisfy that 
$G\leq \GL(4,\bZ)$ and $G^\prime\leq \GL(4,\bZ)$ are 
$\GL(4,\bZ)$-conjugate to $N_{4,i}$ and $N_{4,j}$ respectively. 
For $1\leq i,j\leq 152$ 
except for the cases $i=j=137,139,145,147$, 
the following conditions are equivalent:\\
{\rm (1)} $T_i$ and $T_j^\prime$ are stably birationally $k$-equivalent;\\
{\rm (2)} 
$G\simeq G^\prime$, 
$L_i=L_j^\prime$, 
$T_i\times_k K$ and $T_j^\prime\times_k K$ 
are weak stably birationally $K$-equivalent 
for any $k\subset K\subset L_i$;\\
{\rm (3)} 
$G\simeq G^\prime$, 
$L_i=L_j^\prime$, 
$T_i\times_k K$ and $T_j^\prime\times_k K$ 
are weak stably birationally $K$-equivalent 
for any $k\subset K\subset L_i$ corresponding to 
${\rm WSEC}_r$ $(1\leq r\leq 128)$ as in Table $4$ and Table $5$ 
with $T_i\times_k K$ not stably $K$-rational;\\
{\rm (4)} 
$G\simeq G^\prime$, 
$L_i=L_j^\prime$, 
$T_i\times_k K$ and $T_j^\prime\times_k K$ 
are weak stably birationally $K$-equivalent 
for any $k\subset K\subset L_i$ corresponding to 
${\rm WSEC}_r$ $(1\leq r\leq 128)$ as in Table $4$ and Table $5$ 
with $T_i\times_k K$ not stably $K$-rational 
and $[K:k]=d$ 
$($the ``bold'' cases as in {\rm Table} $12$$)$ 
where 
\begin{align*}
d=
\begin{cases}
1&(i=1,2,4,5,7,10,13,14,23,24,27,39,41,42,44,49,57,58,59,60,61,62,65,67,69,70,71,73,\\
&\hspace*{7mm}74,75,76,77, 78,79,81,82,84,87,91,92,94,95,97,98,101,104,106,107,108,109,110,111,\\
&\hspace*{7mm}114,115,116,118,119,121,124,126,127,130,132,133,134,138
%%%%%%%
,137,139,
%%%%%%%
140,141,
%%%%%%%
145),\\
1,2&(i=3,6,8,9,11,12,15,16,17,18,19,20,22,25,26,29,30,31,33,35,37,43,46,50,51,54,55,\\
&\hspace*{7mm}63,64,66,68,72,80,83,85,89,90,93,96,99,100,102,103,105,125),\\
1,3&(i=142,148),\\
1,4&(i=21,36,38,45,47,48,52,53,56,113,117,120),\\
1,2,4&(i=28,32,34,40),\\
1,6&(i=86,88),\\
1,8&(i=112,122,123,128,129,131,135,136),\\
1,3,8&(i=146),\\
1,12&(i=143,147),\\
1,24&(i=144,149),\\
1,32&(i=150),\\
1,2,32&(i=152),\\
1,3,32&(i=151).
\end{cases}
\end{align*}
In particular, for the cases with $d=1$, we get that 
${(\rm 4)}$ 
if and only if $G\simeq G^\prime$, $L_i=L_j^\prime$, 
i.e. $\widetilde{H}\simeq G\simeq G^\prime$. 

For the exceptional cases $i=j=137,139,145,147$ 
$($$G\simeq Q_8\times C_3$, $(Q_8\times C_3)\rtimes C_2$, 
$\SL(2,\bF_3)\rtimes C_4$, $(\GL(2,\bF_3)\rtimes C_2)\rtimes C_2\simeq 
(\SL(2,\bF_3)\rtimes C_4)\rtimes C_2$$)$, 
we have the implications 
${\rm (1)}\Rightarrow {\rm (2)}\Leftrightarrow {\rm (3)}
\Leftrightarrow {\rm (4)}$ and if $G\simeq G^\prime$, $L_i=L_j^\prime$, then 
there exists $\sigma\in {\rm Aut}(G)$ such that $G^\prime=G^\sigma$ 
and $X=Y\lhd Z$ with $Z/Y\simeq C_2,C_2^2,C_2,C_2$ respectively 
$($this means that the implication $(1)\Rightarrow (2)$ 
can not be reversed in general$)$
where ${\rm Inn}(G)\leq X\leq Y\leq Z\leq {\rm Aut}(G)$ 
are as in Definition \ref{defXYZ} 
%where 
%\begin{center}
%${\rm Inn}(G)\leq X\leq Y\leq Z\leq {\rm Aut}(G)$,
%\end{center}
%\begin{align*}
%X&={\rm Aut}_{\GL(4,\bZ)}(G)=
%\{\sigma\in{\rm Aut}(G)\mid 
%{\rm there}\ {\rm exists}\ u\in {\rm GL}(4,\bZ)\ {\rm such}\ {\rm that}\ 
%u^{-1}gu=g^\sigma\ {\rm for}\ {\rm any}\ g\in G\}\\
%&\simeq N_{\GL(4,\bZ)}(G)/Z_{\GL(4,\bZ)}(G),\\
%Y&=\{\sigma\in{\rm Aut}(G)\mid [M_G]^{fl}=[M_{G^\sigma}]^{fl}\ {\rm as}\ 
%\widetilde{H}\textrm{-{\rm lattices}}\ {\rm where}\ \widetilde{H}=\{(g,g^\sigma)\mid g\in G\}\simeq G\},\\
%Z&=\{\sigma\in{\rm Aut}(G)\mid [M_H]^{fl}\sim [M_{H^\sigma}]^{fl}\ {\rm for}\ {%\rm any}\ H\leq G\}, 
%\end{align*}
%${\rm Inn}(G)$ is the group of inner automorphisms on $G$, 
%${\rm Aut}(G)$ is the group of automorphisms on $G$, 
%$N_{\GL(4,\bZ)}(G)$ is the normalizer of $G$ in $\GL(4,\bZ)$ and 
%$Z_{\GL(4,\bZ)}(G)$ is the centralizer of $G$ in $\GL(4,\bZ)$ 
%
and we have 
${\rm (1)}$ 
$\Leftrightarrow M_G\simeq M_{G^\sigma}$ as $\widetilde{H}$-lattices 
$\Leftrightarrow M_G\otimes_\bZ \bF_p\simeq M_{G^\sigma}\otimes_\bZ \bF_p$ 
as $\bF_p[\widetilde{H}]$-lattices for
\begin{align*}
p= 
\begin{cases}
2&(i=j=137),\\
2\ \,{\rm and}\ \,3&(i=j=139),\\
3&(i=j=145,147).
\end{cases}
\end{align*}

Moreover, for $1\leq r\leq 121$, $r\neq 7,8,10,12$, 
we get the following disjoint union decompositions 
\begin{align*}
{\rm WSEC}_r=\coprod_{L/k\atop {\rm Gal}(L/k)\simeq {\rm G}_r}
{\rm WSEC}_{r,L},\ 
{\rm WSEC}_{r,L}=
\coprod_{t=1}^{\lambda_{r}} {\rm SEC}_{r,L,t}
\end{align*}
modulo stably birationally $k$-equivalence $\stackrel{\rm s.b.}{\approx}$ 
where 
${\rm SEC}_{r,L,t}$ $(1\leq t\leq \lambda_{r})$ is the $t$-th 
stably $k$-equivalent class of $T$ of dimension $4$ in 
${\rm WSEC}_{r,L}$ which corresponds to the fixed minimal splitting field $L$ 
in 
${\rm WSEC}_r=\coprod_{L/k\atop {\rm Gal}(L/k)\simeq {\rm G}_r} {\rm WSEC}_{r,L}$ 
with 
$[\widehat{T}]^{fl}\in {\rm WSEC}_{r,L}
=\coprod_{t=1}^{\lambda_r} {\rm SEC}_{r,L,t}$, 
${\rm Gal}(L/k)\simeq {\rm G}_r\simeq N_{4,i}$ 
$(1\leq r\leq 121, r\neq 7,8,10,12)$ and 
$\lambda_r=|{\rm WSEC}_{r,L}|=|Y\backslash {\rm Aut}(G)|$ 
is given as in Table $4$. 
%
Furthermore, 
for the box cases \fbox{$N_{4,i}$} with $X=Y$ as in Table $4$,
the following conditions are also equivalent:\\
{\rm (0)} $T_i$ and $T_i^\prime$ are birationally $k$-equivalent;\\
{\rm (1)} $T_i$ and $T_i^\prime$ are stably birationally $k$-equivalent. 

In particular, for the box cases \fbox{$N_{4,i}$}, 
all the algebraic tori $T$ of dimension $4$ 
with $\widehat{T}=M_G$ and $G=N_{4,i}$ 
which correspond to the fixed minimal splitting field $L$ 
split into 
$\lambda_r=|Y\backslash {\rm Aut}(G)|$ birationally $k$-equivalent classes. 
\end{theorem}

\begin{theorem}\label{thmXYZ}
Let $k$ be a field. 
Let $G, G^\sigma\leq {\rm GL}(n,\bZ)$ $(\sigma\in {\rm Aut}(G))$ 
and $M_G$, $M_{G^\sigma}$ be the corresponding $G$-lattices 
as in Definition \ref{d2.2}. 
Let $T$ and $T^\sigma$ be algebraic $k$-tori of dimension $n$ 
with the minimal splitting field $L$ and the character modules 
$\widehat{T}=M_G$ and $\widehat{T}^\sigma=M_{G^\sigma}$ 
$(\sigma\in {\rm Aut}(G))$
with 
%$G, G^\sigma\leq \GL(n,\bZ)$, 
$G\simeq G^\sigma\simeq {\rm Gal}(L/k)$. 
Let ${\rm Inn}(G)\leq X\leq Y\leq Z\leq {\rm Aut}(G)$ 
be as in Definition \ref{defXYZ}.\\
{\rm (1)} For $\sigma\in {\rm Aut}(G)$, 
$T$ and $T^\sigma$ are weak stably birationally $k$-equivalent;\\
{\rm (2)} For $\sigma\in {\rm Aut}(G)$, 
$\sigma\in X$ if and only if 
$M_G\simeq M_{G^\sigma}$ as ${\rm Gal}(L/k)$-lattices. 
In particular,  
$T$ and $T^\sigma$ are birationally $k$-equivalent, i.e. 
$k(T)\simeq L(M)^G\simeq L(M^\sigma)^{G^\sigma}\simeq k(T^\sigma)$;\\
{\rm (3)} For $\sigma\in {\rm Aut}(G)$, 
$\sigma\in Y$ if and only if 
$T$ and $T^\sigma$ are stably birationally $k$-equivalent. 
In particular, 
$\{T^\sigma\mid \sigma\in {\rm Aut}(G)\}$ splits into 
$\lambda$ stably birationally $k$-equivalent classes 
where 
$\lambda=|Y\backslash {\rm Aut}(G)|$;\\ %$Y=Z$. 
{\rm (4)} For $\sigma\in {\rm Aut}(G)$, 
$\sigma\in Z$ if and only if 
$T\times_k K$ and $T^\sigma\times_k K$ 
are weak stably birationally $K$-equivalent for any $k\subset K\subset L$. 
%{\rm (3)} $X=Y$ if and only if 
%$M\simeq M^\sigma$ as ${\rm Gal}(L/k)$-lattices for any $\sigma\in Y$;\\
%{\rm (4)} $Y=Z$ if and only if $(2)\Rightarrow (1)$;\\

In particular, {\rm (i)} if $Y=Z$, then 
$\sigma \in Z$ 
if and only if 
$T$ and $T^\sigma$ are stably birationally $k$-equivalent; 
and {\rm (ii)} if $X=Y=Z$ $($resp. $X=Y$$)$, then 
$\sigma\in Z$ $($resp. $\sigma\in Y$$)$
if and only if 
$T$ and $T^\sigma$ are birationally $k$-equivalent. 
\end{theorem}

\begin{proof}
(1) follows from 
$[M_G]^{fl}\sim [M_{G^\sigma}]^{fl}$ 
because 
$[M_G]^{fl}=[M_{G^\sigma}]^{fl}$ 
as $\widetilde{H}$-lattices where 
$\widetilde{H}=\{(g,(g^{\sigma^{-1}})^{\sigma})\mid 
g\in G\}\leq G\times G^\sigma$ and $\widetilde{H}\simeq G\simeq G^\sigma$. 

Because $T$ and $T^\sigma$ have the same splitting field $L$ 
with ${\rm Gal}(L/k)\simeq G\simeq G^\sigma$, 
we have\\
(2) $\sigma\in X$
$\Leftrightarrow$ 
there exists $u\in {\rm GL}(n,\bZ)$ such that 
$u^{-1}gu=g^\sigma$ for any $g\in G$ 
$\Leftrightarrow$ 
$M_G\simeq M_{G^\sigma}$ as ${\rm Gal}(L/k)$-lattices;\\
(3) 
%Because $T$ and $T^\sigma$ have the same splitting field $L$ 
%with ${\rm Gal}(L/k)\simeq G$, 
%we have 
$\sigma\in Y$
$\Leftrightarrow$ 
$[M_G]^{fl}=[M_{G^\sigma}]^{fl}$ as 
$\widetilde{H}$-lattices where 
$\widetilde{H}=\{(g,g^\sigma)\mid g\in G\}\simeq G$ 
$\Leftrightarrow$ 
$T$ and $T^\sigma$ are stably birationally $k$-equivalent;\\
(4) 
$\sigma\in Z$
$\Leftrightarrow$ 
$[M_H]^{fl}\sim [M_{H^\sigma}]^{fl}$ for any 
$H\leq G$ 
$\Leftrightarrow$ 
$T\times_k K$ and $T^\sigma\times_k K$ 
are weak stably birationally $K$-equivalent 
for any $k\subset K\subset L$. 
\end{proof}

\end{document}
